\chapter{A Taste of Modular Forms}

\begin{center}\itshape
There are five elementary arithmetical operations: addition, subtraction, multiplication, division, and modular forms.\\
\upshape --- attributed to Martin Eichler
\end{center}

In 1916 Srinivasa Ramanujan, two years after arriving in Cambridge, published a paper on sums of divisors. In it he expanded the infinite product
\[
\begin{aligned}q\prod_{n\ge1}(1-q^n)^{24}&=\sum_{n\ge1}\tau(n)q^n\\&=q-24q^2+252q^3-1472q^4\eqbr\quad+4830q^5-\cdots\end{aligned}
\]
and looked at the coefficients. They seemed to have no reason to be anything in particular. Ramanujan noticed three things. First, $\tau(6)=-6048=\tau(2)\tau(3)$, and more generally $\tau(mn)=\tau(m)\tau(n)$ whenever
$\gcd(m,n)=1$. Second, $\tau(n)\equiv\sigma_{11}(n)\pmod{691}$, where $\sigma_{11}(n)$ is the sum of the eleventh powers of the divisors of $n$. Third, $|\tau(p)|\le2p^{11/2}$ for every prime $p$. He could prove none of them.

Mordell proved the first a year later. The second follows from a single identity between two functions, and we shall prove it in a few lines. The third resisted for fifty-eight years: Deligne proved it in 1974,
as a consequence of his proof of the Weil conjectures, one of the deepest results of twentieth-century mathematics. All three are statements about \emph{modular forms}, functions on the upper half-plane with an
enormous group of symmetries. That group forces the space of such functions to be finite-dimensional, and finite-dimensional spaces are where coincidences become theorems.

\begin{think}
Two power series built from divisor sums, $1+240\sum\sigma_3(n)q^n$ and $1+480\sum\sigma_7(n)q^n$, turn out to satisfy $(1+240\sum\sigma_3(n)q^n)^2=1+480\sum\sigma_7(n)q^n$. Compare the coefficients of $q^2$ by hand.
What identity between divisor sums does this give, and why would anyone expect it to hold for every $n$?
\end{think}

This last chapter is a taste, not a meal. We prove what is needed to explain Ramanujan's first two observations, with full proofs, and use only complex analysis from Tome~4. Tome~12 develops the theory properly.

\section{The modular group}

Let $\mathbb H=\{\tau\in\C:\operatorname{Im}\tau>0\}$ be the upper half-plane. The group $\mathrm{SL}_2(\Z)$ of integer matrices of determinant $1$ acts on $\mathbb H$ by Möbius transformations:
\[
\gamma=\begin{pmatrix}a&b\\c&d\end{pmatrix},\qq\gamma\tau=\frac{a\tau+b}{c\tau+d},\qq\operatorname{Im}(\gamma\tau)=\frac{\operatorname{Im}\tau}{|c\tau+d|^2}.
\]
The last formula follows from $(a\tau+b)(c\overline\tau+d)=ac|\tau|^2+bd+ad\tau+bc\overline\tau$, whose imaginary part is $(ad-bc)\operatorname{Im}\tau$. So $\gamma\tau\in\mathbb H$, and one checks that $(\gamma\gamma')\tau=\gamma(\gamma'\tau)$
by multiplying out. The matrices $\gamma$ and $-\gamma$ act in the same way. Two elements play a special role:
\[
\begin{gathered}T=\begin{pmatrix}1&1\\0&1\end{pmatrix},\qquad T\tau=\tau+1,\\S=\begin{pmatrix}0&-1\\1&0\end{pmatrix},\qquad S\tau=-\frac1\tau.\end{gathered}
\]

\begin{proposition}\label{prop:generators}
The matrices $S$ and $T$ generate $\mathrm{SL}_2(\Z)$.
\end{proposition}

\begin{proof}
Let $G$ be the subgroup generated by $S$ and $T$; it contains $S^2=-I$. Let $\gamma=\bigl(\begin{smallmatrix}a&b\\c&d\end{smallmatrix}\bigr)\in\mathrm{SL}_2(\Z)$; we show $\gamma\in G$ by induction on $|c|$. If $c=0$, then $ad=1$, so $\gamma=\pm\bigl(\begin{smallmatrix}1&b'\\0&1\end{smallmatrix}\bigr)=\pm T^{b'}$.
If $c\ne0$, Euclid's division gives $d=kc+r$ with $0\le r<|c|$, and
\[
\gamma T^{-k}=\begin{pmatrix}a&b-ka\\c&r\end{pmatrix},\qq\gamma T^{-k}S=\begin{pmatrix}b-ka&-a\\r&-c\end{pmatrix}.
\]
The lower-left entry $r$ satisfies $|r|<|c|$, so $\gamma T^{-k}S\in G$ by induction, and so is $\gamma$.
\end{proof}

\begin{theorem}\label{thm:fundamental}
Let $F=\{\tau\in\mathbb H:|\operatorname{Re}\tau|\le\frac12,\ |\tau|\ge1\}$. Every $\tau\in\mathbb H$ is of the form $\gamma\tau'$ with $\gamma\in\mathrm{SL}_2(\Z)$ and $\tau'\in F$.
\end{theorem}

\begin{proof}
Fix $\tau$. For $c,d$ integers, $|c\tau+d|\le1$ forces $|c|\operatorname{Im}\tau\le1$ and then $|d|\le1+|c||\tau|$, so only finitely many pairs $(c,d)$ satisfy it. Hence $\operatorname{Im}(\gamma\tau)=\frac{\operatorname{Im}\tau}{|c\tau+d|^2}$
takes only finitely many values larger than $\operatorname{Im}\tau$, and there is $\gamma$ for which $\operatorname{Im}(\gamma\tau)$ is maximal. Replacing $\gamma$ by $T^m\gamma$, which does not change the imaginary part, we may assume $|\operatorname{Re}\gamma\tau|\le\frac12$.
If $|\gamma\tau|<1$, then $\operatorname{Im}(S\gamma\tau)=\frac{\operatorname{Im}\gamma\tau}{|\gamma\tau|^2}>\operatorname{Im}\gamma\tau$, contradicting maximality. So $\gamma\tau\in F$, and $\tau=\gamma^{-1}(\gamma\tau)$.
\end{proof}

So $F$ meets every orbit, and Tome~12 shows that it meets each orbit only once, apart from identifications on its boundary. The images of $F$ tile $\mathbb H$ (Figure~\ref{fig:tessellation}). Every point of $F$
has imaginary part at least $\frac{\sqrt3}2$, the height of the corner $\rho=e^{2\pi i/3}$.

\begin{figure}[ht]
\centering
\fig{ch12-tessellation}
\caption{The fundamental domain $F$ (shaded) and its images under $\mathrm{SL}_2(\Z)$, which tile the upper half-plane. Each region is bounded by arcs of circles centred on the real axis, or vertical lines. The tiles
accumulate at every rational point of the real axis.}
\label{fig:tessellation}
\end{figure}

\section{Modular forms and Eisenstein series}

For $\gamma=\bigl(\begin{smallmatrix}a&b\\c&d\end{smallmatrix}\bigr)$ put $j(\gamma,\tau)=c\tau+d$. Multiplying out shows the \emph{cocycle relation}
\[
j(\gamma\gamma',\tau)=j(\gamma,\gamma'\tau)\,j(\gamma',\tau).
\]
For an integer $k$ and a function $f$ on $\mathbb H$, define $(f|_k\gamma)(\tau)=j(\gamma,\tau)^{-k}f(\gamma\tau)$. The cocycle relation says exactly that $f|_k(\gamma\gamma')=(f|_k\gamma)|_k\gamma'$.

\begin{definition}\label{def:modform}
A \emph{modular form of weight $k$} is a holomorphic function $f$ on $\mathbb H$ such that $f|_k\gamma=f$ for all $\gamma\in\mathrm{SL}_2(\Z)$, that is
\[
f\Bigl(\frac{a\tau+b}{c\tau+d}\Bigr)=(c\tau+d)^kf(\tau),
\]
and which is \emph{holomorphic at infinity}: there is a power series $h(q)=\sum_{n\ge0}a_nq^n$ converging for $|q|<1$ with $f(\tau)=h(e^{2\pi i\tau})$. The $a_n$ are the \emph{Fourier coefficients} of $f$; if $a_0=0$, $f$ is a
\emph{cusp form}. We write $M_k$ and $S_k$ for the complex vector spaces of modular forms and of cusp forms of weight $k$.
\end{definition}

From now on $q=e^{2\pi i\tau}$; note $|q|=e^{-2\pi\operatorname{Im}\tau}<1$. By Proposition~\ref{prop:generators} and the composition rule, $f|_k\gamma=f$ holds for all $\gamma$ as soon as it holds for $S$ and $T$, that is,
as soon as $f(\tau+1)=f(\tau)$ and $f(-\frac1\tau)=\tau^kf(\tau)$. Taking $\gamma=-I$ gives $f=(-1)^kf$, so $M_k=0$ for odd $k$. Products of modular forms are modular, and weights add: $M_kM_l\subseteq M_{k+l}$.

The first examples are lattice sums. For an even $k\ge4$ and $\tau\in\mathbb H$, the \emph{Eisenstein series} is
\[
G_k(\tau)=\sum_{(m,n)\ne(0,0)}\frac1{(m\tau+n)^k}.
\]

\begin{lemma}\label{lem:latticesum}
For $\sigma>2$ the series $\sum_{(m,n)\ne(0,0)}|m\tau+n|^{-\sigma}$ converges uniformly for $\tau$ in every compact subset of $\mathbb H$.
\end{lemma}

\begin{proof}
Let $K\subseteq\mathbb H$ be compact. The function $(\tau,x,y)\mapsto|x\tau+y|$ is continuous and positive on the compact set of $(\tau,x,y)$ with $\tau\in K$ and $\max(|x|,|y|)=1$ (it vanishes only if $x=y=0$, since $\tau\notin\R$),
so it has a minimum $c>0$ there. By homogeneity, $|m\tau+n|\ge c\max(|m|,|n|)$ for all $(m,n)$. There are $8r$ pairs with $\max(|m|,|n|)=r\ge1$, so the series is at most $\sum_{r\ge1}8r(cr)^{-\sigma}<\infty$, uniformly on $K$.
\end{proof}

\begin{proposition}\label{prop:eisenstein}
For even $k\ge4$, $G_k$ is holomorphic on $\mathbb H$ and $G_k|_k\gamma=G_k$ for all $\gamma\in\mathrm{SL}_2(\Z)$.
\end{proposition}

\begin{proof}
Holomorphy follows from Lemma~\ref{lem:latticesum} and Weierstrass's theorem (Tome~4, Theorem~9.15). For $\gamma=\bigl(\begin{smallmatrix}a&b\\c&d\end{smallmatrix}\bigr)$,
\[
\frac1{(m\gamma\tau+n)^k}=\frac{(c\tau+d)^k}{\bigl((ma+nc)\tau+(mb+nd)\bigr)^k},
\]
and $(m,n)\mapsto(ma+nc,mb+nd)$ is a bijection of $\Z^2\setminus\{0\}$, with inverse given by the inverse matrix. Summing over $(m,n)$, which is allowed by absolute convergence, $G_k(\gamma\tau)=(c\tau+d)^kG_k(\tau)$.
\end{proof}

\begin{theorem}[Fourier expansion of $G_k$]\label{thm:eisensteinq}
For even $k\ge4$,
\[
G_k(\tau)=2\zeta(k)+\frac{2(2\pi i)^k}{(k-1)!}\sum_{n\ge1}\sigma_{k-1}(n)q^n,
\]
where $\sigma_{k-1}(n)=\sum_{d\mid n}d^{k-1}$. Consequently $E_k=\frac{G_k}{2\zeta(k)}$ is a modular form of weight $k$, with
\[
E_k=1-\frac{2k}{B_k}\sum_{n\ge1}\sigma_{k-1}(n)q^n.
\]
\end{theorem}

\begin{proof}
By Tome~4, Corollary~9.35, for $\tau\in\mathbb H$,
\[
\pi\cot\pi\tau=\frac1\tau+\sum_{n\ge1}\Bigl(\frac1{\tau+n}+\frac1{\tau-n}\Bigr),
\]
and on the other hand $\pi\cot\pi\tau=\pi i\frac{q+1}{q-1}=-\pi i\bigl(1+2\sum_{r\ge1}q^r\bigr)$ since $|q|<1$. Both series converge uniformly on compact subsets of $\mathbb H$, so we may differentiate $k-1$ times term by term
(Tome~4, Theorem~9.15). As $\frac{\mathrm d^{k-1}}{\mathrm d\tau^{k-1}}\frac1{\tau+n}=\frac{(-1)^{k-1}(k-1)!}{(\tau+n)^k}$ and $\frac{\mathrm d}{\mathrm d\tau}q^r=2\pi irq^r$, we get, for even $k$, the \emph{Lipschitz formula}
\[
\sum_{n\in\Z}\frac1{(\tau+n)^k}=\frac{(2\pi i)^k}{(k-1)!}\sum_{r\ge1}r^{k-1}q^r.
\]
In $G_k$, the terms with $m=0$ give $2\zeta(k)$. The terms with $m\ne0$ come in pairs $(m,n)$, $(-m,-n)$ with equal values since $k$ is even, so they give $2\sum_{m\ge1}\sum_n(m\tau+n)^{-k}$. Apply the Lipschitz formula
at $m\tau$, whose $q$ is $q^m$:
\[
\begin{aligned}G_k(\tau)-2\zeta(k)&=\frac{2(2\pi i)^k}{(k-1)!}\sum_{m\ge1}\sum_{r\ge1}r^{k-1}q^{mr}\eqbr=\frac{2(2\pi i)^k}{(k-1)!}\sum_{n\ge1}\sigma_{k-1}(n)q^n,\end{aligned}
\]
grouping the terms with $mr=n$ (the double series converges absolutely). Finally, Euler's formula (Tome~4, Theorem~3.7) gives $2\zeta(k)=\frac{(-1)^{k/2+1}B_k(2\pi)^k}{k!}$, and $(2\pi i)^k=(-1)^{k/2}(2\pi)^k$, so
$\frac{2(2\pi i)^k}{(k-1)!\,2\zeta(k)}=-\frac{2k}{B_k}$.
\end{proof}

With $B_4=-\frac1{30}$, $B_6=\frac1{42}$, $B_8=-\frac1{30}$, $B_{10}=\frac5{66}$, $B_{12}=-\frac{691}{2730}$:
\[
\begin{gathered}E_4=1+240\sum\sigma_3(n)q^n,\qq E_6=1-504\sum\sigma_5(n)q^n,\\E_8=1+480\sum\sigma_7(n)q^n,\qq E_{10}=1-264\sum\sigma_9(n)q^n,\\E_{12}=1+\frac{65520}{691}\sum\sigma_{11}(n)q^n.\end{gathered}
\]
There it is: the prime $691$, hiding in the numerator of a Bernoulli number.

\section{Dedekind's eta function}

Ramanujan's product is the $24$-th power of a function that Dedekind studied in 1877:
\[
\begin{gathered}\eta(\tau)=e^{\pi i\tau/12}\prod_{n\ge1}(1-q^n),\\\Delta(\tau)=\eta(\tau)^{24}=q\prod_{n\ge1}(1-q^n)^{24}.\end{gathered}
\]
Since $\sum|q|^n$ converges, uniformly on $\operatorname{Im}\tau\ge y_0>0$, the function $L(\tau)=\frac{\pi i\tau}{12}+\sum_{n\ge1}\log(1-q^n)$ (principal logarithms, as $|q^n|<1$) is holomorphic on $\mathbb H$, and $\eta=e^L$. In particular $\eta$ and
$\Delta$ are holomorphic and \emph{never zero} on $\mathbb H$. Clearly $\Delta(\tau+1)=\Delta(\tau)$. The heart of the matter is the behaviour under $S$.

\begin{theorem}[Dedekind]\label{thm:eta}
For $\tau\in\mathbb H$, $\eta(-\frac1\tau)=\sqrt{-i\tau}\,\eta(\tau)$, with the principal square root. Hence $\Delta(-\frac1\tau)=\tau^{12}\Delta(\tau)$, and $\Delta$ is a cusp form of weight $12$ with $\tau(1)=1$.
\end{theorem}

The proof below is Siegel's (1954): a single contour integral of cotangents.

\begin{proof}
Both $L(-\frac1\tau)$ and $L(\tau)+\frac12\log(-i\tau)$ are holomorphic on $\mathbb H$ ($-i\tau$ lies in the right half-plane), so by the identity theorem (Tome~4, Theorem~9.18) it suffices to prove that they agree for $\tau=iy$, $y>0$.
Expanding $\log(1-x)=-\sum_m\frac{x^m}m$ and summing the geometric series in $n$,
\[
L(iy)=-\frac{\pi y}{12}-\Sigma(y),\qq\Sigma(y)=\sum_{m\ge1}\frac1m\cdot\frac1{e^{2\pi my}-1}.
\]
Since $-\frac1{iy}=\frac iy$, we must show
\[
\Sigma(y)-\Sigma\Bigl(\frac1y\Bigr)+\frac\pi{12}\Bigl(y-\frac1y\Bigr)=\frac12\log y.\tag{$\ast$}
\]
Let $\tau=iy$, $N=n+\frac12$ with $n\ge1$ an integer, and
\[
f_N(z)=\frac1z\cot(\pi Nz)\cot\Bigl(\frac{\pi Nz}\tau\Bigr).
\]
Let $C$ be the parallelogram (here a rhombus) with vertices $1,\tau,-1,-\tau$, traversed positively. Inside $C$, $f_N$ has poles at $z=0$, at $z=\frac kN$ and at $z=\frac{k\tau}N$ for $1\le|k|\le n$; none lies on $C$, since
$N$ is not an integer.

\emph{Residues.} At $z=\frac kN$, $\cot(\pi Nz)$ has residue $\frac1{\pi N}$, so $f_N$ has residue $\frac1{\pi k}\cot\frac{\pi k}\tau$. At $z=\frac{k\tau}N$, similarly, the residue is $\frac1{\pi k}\cot(\pi k\tau)$. At $z=0$, from
$\cot u=\frac1u-\frac u3+O(u^3)$,
\[
\begin{aligned}\cot(\pi Nz)\cot\Bigl(\frac{\pi Nz}\tau\Bigr)&=\frac\tau{\pi^2N^2z^2}\eqbr-\frac13\Bigl(\tau+\frac1\tau\Bigr)+O(z^2),\end{aligned}
\]
so the residue is $-\frac13(\tau+\frac1\tau)$. The terms $k$ and $-k$ give equal residues. Now $\cot(ix)=-i\coth x=-i\bigl(1+\frac2{e^{2x}-1}\bigr)$, so $\cot(\pi k\tau)=-i\bigl(1+\frac2{e^{2\pi ky}-1}\bigr)$ and
$\cot\frac{\pi k}\tau=\cot\frac{-i\pi k}y=i\bigl(1+\frac2{e^{2\pi k/y}-1}\bigr)$; in the sum the constants cancel. With $\tau+\frac1\tau=i(y-\frac1y)$, the residue theorem (Tome~4, Theorem~9.26) gives
\[
\begin{aligned}\frac1{2\pi i}\int_Cf_N&=-\frac i3\Bigl(y-\frac1y\Bigr)\eqbr+\frac{4i}\pi\sum_{k=1}^n\frac1k\Bigl(\frac1{e^{2\pi k/y}-1}-\frac1{e^{2\pi ky}-1}\Bigr).\end{aligned}
\]
Multiplying by $\frac\pi{4i}$: $-\frac18\int_Cf_N=-\frac\pi{12}(y-\frac1y)+\Sigma_n(\frac1y)-\Sigma_n(y)$, where $\Sigma_n$ is the partial sum up to $n$.

\emph{The limit $N\to\infty$.} For $w=u+iv$, $|\cot w|$ is bounded on the set of $w$ at distance at least $\delta$ from $\pi\Z$, for each $\delta>0$, and $\cot w\to\mp i$ as $v\to\pm\infty$. On the edge from $1$ to $\tau$,
$z=1-t+ity$ with $0\le t\le1$, and $\pi Nz=\pi(n+\frac12)-\pi Nt+i\pi Nty$; its distance to the poles $\pi m$ is at least $\min_{s\ge0}|\frac\pi2-s+isy|=\frac\pi2\frac y{\sqrt{1+y^2}}$, independently of $N$. The same holds for $\frac{\pi Nz}\tau$ and on the other
edges, by symmetry. So $z f_N(z)$ is bounded on $C$ uniformly in $N$. At each point of the open edges, $\operatorname{Im}(\pi Nz)$ and $\operatorname{Im}(\frac{\pi Nz}\tau)=-\frac{\pi N\operatorname{Re}z}y$ tend to $\pm\infty$, and $zf_N(z)$ tends to
$+1$ on the edges $[1,\tau]$ and $[-1,-\tau]$, where $\operatorname{Re}z$ and $\operatorname{Im}z$ have the same sign, and to $-1$ on the two others. By bounded convergence,
\[
\begin{aligned}\int_Cf_N\to&\int_{[1,\tau]}\frac{\mathrm dz}z-\int_{[\tau,-1]}\frac{\mathrm dz}z\eqbr+\int_{[-1,-\tau]}\frac{\mathrm dz}z-\int_{[-\tau,1]}\frac{\mathrm dz}z.\end{aligned}
\]
Along each edge, $\int\frac{\mathrm dz}z$ is the change of $\log|z|+i\arg z$, with the argument varying continuously: the four integrals are $\log y+\frac{i\pi}2$, $-\log y+\frac{i\pi}2$, $\log y+\frac{i\pi}2$, $-\log y+\frac{i\pi}2$. So the limit is $4\log y$, and
letting $n\to\infty$ in the identity above gives $-\frac12\log y=-\frac\pi{12}(y-\frac1y)+\Sigma(\frac1y)-\Sigma(y)$, which is $(\ast)$.

Exponentiating, $\eta(-\frac1\tau)=e^{\frac12\log(-i\tau)}\eta(\tau)$, and $\Delta(-\frac1\tau)=(-i\tau)^{12}\Delta(\tau)=\tau^{12}\Delta(\tau)$. With $\Delta(\tau+1)=\Delta(\tau)$, Proposition~\ref{prop:generators} gives $\Delta|_{12}\gamma=\Delta$ for all $\gamma$.
Expanding the product gives a power series $\sum_{n\ge1}\tau(n)q^n$ with integer coefficients and $\tau(1)=1$.
\end{proof}

\section{Finite dimensions}

Now the symmetry does its work. A modular form is determined by its values on $F$, and $F$ is ``compact up to the point at infinity''. This is enough to make weight-$0$ forms constant.

\begin{lemma}\label{lem:weight0}
Every modular form of weight $0$ is constant.
\end{lemma}

\begin{proof}
Let $f\in M_0$, $f(\tau)=h(q)$ with $h$ holomorphic on the disc $|q|<1$. Since $f$ is invariant, Theorem~\ref{thm:fundamental} shows that $f$ takes on $\mathbb H$ the same values as on $F$. Points of $F$ have
$\operatorname{Im}\tau\ge\frac{\sqrt3}2$, hence $|q|\le r_0=e^{-\pi\sqrt3}$. So
\[
\sup_{|q|<1}|h|=\sup_F|f|\le\max_{|q|\le r_0}|h|,
\]
and $|h|$ attains its supremum over the open disc at some point $q_0$ with $|q_0|\le r_0$. By the
maximum principle, $h$ is constant. (Proof of the principle: by Cauchy's formula $h(q_0)$ is the mean of $h$ on each small circle around $q_0$, so if $|h(q_0)|$ is maximal, $|h|$ is constant near $q_0$; then $\overline h=\frac{|h|^2}h$ is
holomorphic there too (or $h=0$), and the Cauchy--Riemann equations, Tome~4, Proposition~9.2, force $h'=0$ near $q_0$; by the identity theorem $h$ is constant.)
\end{proof}

\begin{theorem}\label{thm:dimension}
\begin{enumerate}[label=(\alph*)]
\item $M_k=0$ for $k<0$ and for $k=2$; $M_0=\C$.
\item For even $k\ge4$, $M_k=\C E_k\oplus S_k$, and multiplication by $\Delta$ is an isomorphism $M_{k-12}\to S_k$.
\item Hence $\dim M_k=\lfloor\frac k{12}\rfloor$ if $k\equiv2\pmod{12}$, and $\lfloor\frac k{12}\rfloor+1$ otherwise, for even $k\ge0$.
\end{enumerate}
\end{theorem}

\begin{proof}
\emph{(b) first.} For $f\in M_k$, $f-a_0(f)E_k$ has constant term $0$, so $M_k=\C E_k+S_k$, and the sum is direct since $E_k\notin S_k$. If $f\in S_k$, then $f/\Delta$ is holomorphic on $\mathbb H$ ($\Delta$ has no zeros), satisfies
$(f/\Delta)|_{k-12}\gamma=f/\Delta$, and is holomorphic at infinity: $h_f(q)=q\,g(q)$ and $h_\Delta(q)=q\,u(q)$ with $g,u$ holomorphic on the disc and $u$ without zeros (as $\Delta\ne0$, and $u(0)=1$). So $f/\Delta\in M_{k-12}$, and
$g\mapsto g\Delta$ is the inverse isomorphism.

\emph{(a).} $M_0=\C$ by Lemma~\ref{lem:weight0}. Let $f\in M_k$ with $k<0$. Then $f^{12}\Delta^{-k}$ has weight $12k-12k=0$, so it is constant; and it vanishes at infinity, since $\Delta$ does and $-k>0$. Hence it is $0$ and
$f=0$. So $M_k=0$ for $k<0$, and by (b), $S_k=0$ for $4\le k\le10$; thus $M_4=\C E_4$ and $M_6=\C E_6$. Let $f\in M_2$. A cusp form of weight $2$ would give $f/\Delta\in M_{-10}=0$, so if $f\ne0$ then $c=a_0(f)\ne0$. Now
$f^2\in M_4$ and $f^3\in M_6$, with constant terms $c^2$ and $c^3$, so $f^2=c^2E_4$ and $f^3=c^3E_6$. Then $c^6E_4^3=f^6=c^6E_6^2$, so $E_4^3=E_6^2$. But $E_4^3-E_6^2=(1+720q+\cdots)-(1-1008q+\cdots)=1728q+\cdots\ne0$.
So $M_2=0$.

\emph{(c)} follows from (a) and (b) by induction on $k$, since $\dim M_k=1+\dim M_{k-12}$ for $k\ge4$, with the values $1,0,1,1,1,1$ for $k=0,2,4,6,8,10$.
\end{proof}

The payoff is immediate. $M_8$, $M_{10}$ and $M_{14}$ are one-dimensional, spanned by $E_8$, $E_{10}$ and $E_{14}$, and they also contain $E_4^2$, $E_4E_6$ and $E_4^2E_6$, with constant term $1$. So $E_4^2=E_8$, $E_4E_6=E_{10}$, and
so on. And $S_{12}$ is the line spanned by $\Delta$.

\begin{corollary}\label{cor:sigma7}
For every $n\ge1$,
\[
\sigma_7(n)=\sigma_3(n)+120\sum_{m=1}^{n-1}\sigma_3(m)\sigma_3(n-m).
\]
\end{corollary}

\begin{proof}
Compare the coefficients of $q^n$ in $E_4^2=E_8$:
\[
\begin{aligned}480\sigma_7(n)&=2\cdot240\sigma_3(n)\eqbr+240^2\sum_{m=1}^{n-1}\sigma_3(m)\sigma_3(n-m),\end{aligned}
\]
and divide by $480$.
\end{proof}

Elementary proofs exist, but they are long; here the identity falls out of a dimension count. For $n=2$ it says $129=9+120$.

\begin{corollary}\label{cor:delta}
$1728\Delta=E_4^3-E_6^2$, and Ramanujan's congruence holds: $\tau(n)\equiv\sigma_{11}(n)\pmod{691}$ for every $n\ge1$.
\end{corollary}

\begin{proof}
$E_4^3-E_6^2$ is a cusp form of weight $12$ with $q$-coefficient $720+1008=1728$, so it is $1728\Delta$, as $S_{12}=\C\Delta$ and $\tau(1)=1$. Next, $E_4^3-E_{12}\in S_{12}$ has $q$-coefficient $720-\frac{65520}{691}=\frac{432000}{691}$, so
\[
691E_4^3=691E_{12}+432000\,\Delta.
\]
The coefficient of $q^n$ gives $691a_n=65520\,\sigma_{11}(n)+432000\,\tau(n)$, where $a_n$, the coefficient of $E_4^3$, is an integer. Modulo $691$, $65520\equiv566\equiv-125$ and $432000\equiv125$, so
$125(\tau(n)-\sigma_{11}(n))\equiv0$, and $125$ is invertible modulo the prime $691$.
\end{proof}

\section{Hecke operators}

Why is $\tau$ multiplicative? In 1937 Hecke found the reason: a family of linear operators on each $M_k$ that commute with each other and, on the one-dimensional space $S_{12}$, can only act by scalars. Let $p$ be a prime.
For $f=\sum a_nq^n\in M_k$, define
\[
\begin{aligned}T_pf(\tau)&=p^{k-1}f(p\tau)+\frac1p\sum_{j=0}^{p-1}f\Bigl(\frac{\tau+j}p\Bigr)\eqbr=\sum_{n\ge0}\bigl(a_{pn}+p^{k-1}a_{n/p}\bigr)q^n,\end{aligned}
\]
with $a_{n/p}=0$ if $p\nmid n$. The second expression follows from $\frac1p\sum_je^{2\pi in(\tau+j)/p}=q^{n/p}$ if $p\mid n$ and $0$ otherwise (a sum of $p$-th roots of unity).

To show that $T_pf$ is modular, extend the slash operator to integer matrices $\alpha=\bigl(\begin{smallmatrix}a&b\\c&d\end{smallmatrix}\bigr)$ with $\det\alpha>0$: $(f|_k\alpha)(\tau)=(\det\alpha)^{k-1}(c\tau+d)^{-k}f(\alpha\tau)$. The cocycle relation holds for these
matrices too (it is an identity between polynomials), and $\det$ is multiplicative, so $f|_k(\alpha\beta)=(f|_k\alpha)|_k\beta$. Let $R_p$ be the set of $p+1$ matrices
\[
\begin{pmatrix}p&0\\0&1\end{pmatrix},\qq\begin{pmatrix}1&j\\0&p\end{pmatrix}\ (0\le j<p).
\]
Then $T_pf=\sum_{\rho\in R_p}f|_k\rho$, as one checks directly.

\begin{lemma}\label{lem:hermite}
Every integer matrix $\alpha$ with $\det\alpha=p$ can be written as $\alpha=\gamma\rho$ with $\gamma\in\mathrm{SL}_2(\Z)$ and $\rho\in R_p$, and $\rho$ is unique.
\end{lemma}

\begin{proof}
\emph{Existence.} Let $\alpha=\bigl(\begin{smallmatrix}a&b\\c&d\end{smallmatrix}\bigr)$ and $g=\gcd(a,c)>0$. By Bézout there are $x,y$ with $xa+yc=g$, and $\gamma_1=\bigl(\begin{smallmatrix}x&y\\-c/g&a/g\end{smallmatrix}\bigr)\in\mathrm{SL}_2(\Z)$ satisfies
$\gamma_1\alpha=\bigl(\begin{smallmatrix}g&b'\\0&d'\end{smallmatrix}\bigr)$, with $gd'=p$, so $d'>0$. Multiplying on the left by $T^{-m}$ replaces $b'$ by $b'-md'$, so we may take $0\le b'<d'$. As $p$ is prime, either $g=p$, $d'=1$, $b'=0$,
or $g=1$, $d'=p$, $0\le b'<p$: the result is in $R_p$.

\emph{Uniqueness.} If $\gamma\rho=\gamma'\rho'$ with $\rho,\rho'\in R_p$, then $\rho'\rho^{-1}=\gamma'^{-1}\gamma\in\mathrm{SL}_2(\Z)$. Both are upper triangular, so $\rho'\rho^{-1}$ is upper triangular with positive rational diagonal entries $\frac{g'}g$, $\frac{d'}d$, which
are integers of product $1$; so $g=g'$ and $d=d'$. Its upper right entry is $\frac{b'-b}d$, an integer, so $b\equiv b'\pmod d$, and $b=b'$.
\end{proof}

\begin{theorem}[Hecke]\label{thm:hecke}
If $f\in M_k$, then $T_pf\in M_k$, and $T_p$ maps $S_k$ to $S_k$.
\end{theorem}

\begin{proof}
$T_pf$ is holomorphic on $\mathbb H$, and holomorphic at infinity by its $q$-expansion, whose constant term is $(1+p^{k-1})a_0$; so cusp forms go to cusp forms. Let $\gamma\in\mathrm{SL}_2(\Z)$. For $\rho\in R_p$, $\det(\rho\gamma)=p$, so by
Lemma~\ref{lem:hermite} $\rho\gamma=\gamma_\rho\rho'$ for some $\gamma_\rho\in\mathrm{SL}_2(\Z)$ and a unique $\rho'=\pi(\rho)\in R_p$. The map $\pi$ is injective: if $\pi(\rho_1)=\pi(\rho_2)=\rho'$, then $\rho_1=\gamma_{\rho_1}\rho'\gamma^{-1}$ and
$\rho_2=\gamma_{\rho_2}\rho'\gamma^{-1}$ give $\rho_1=(\gamma_{\rho_1}\gamma_{\rho_2}^{-1})\rho_2$, and uniqueness in the lemma (applied to $\alpha=\rho_1$) forces $\rho_1=\rho_2$. So $\pi$ is a permutation of $R_p$, and
\[
\begin{aligned}(T_pf)|_k\gamma&=\sum_\rho f|_k\rho\gamma=\sum_\rho(f|_k\gamma_\rho)|_k\pi(\rho)\eqbr=\sum_\rho f|_k\pi(\rho)=T_pf.\end{aligned}
\]
\end{proof}

\begin{theorem}[Mordell, 1917]\label{thm:mordell}
For every prime $p$ and $n\ge1$, $\tau(p)\tau(n)=\tau(pn)+p^{11}\tau(n/p)$, with $\tau(n/p)=0$ if $p\nmid n$. Consequently $\tau(mn)=\tau(m)\tau(n)$ whenever $\gcd(m,n)=1$, and
$\tau(p^{r+1})=\tau(p)\tau(p^r)-p^{11}\tau(p^{r-1})$ for $r\ge1$.
\end{theorem}

\begin{proof}
By Theorem~\ref{thm:hecke}, $T_p\Delta\in S_{12}=\C\Delta$, so $T_p\Delta=\lambda\Delta$, and comparing the coefficients of $q$, $\lambda=\tau(p)+p^{11}\tau(1/p)=\tau(p)$. Comparing the coefficients of $q^n$ gives the first identity.
With $n=p^r$ it is the recursion. With $p\nmid n$ it gives $\tau(pn)=\tau(p)\tau(n)$; more generally, for $p\nmid n$ the sequences $r\mapsto\tau(p^rn)$ and $r\mapsto\tau(p^r)\tau(n)$ satisfy the same recursion (apply the first identity to $p^rn$)
and agree for $r=0,1$, so $\tau(p^rn)=\tau(p^r)\tau(n)$. Induction on the number of prime factors of $m$ gives $\tau(mn)=\tau(m)\tau(n)$ for coprime $m,n$.
\end{proof}

So $\tau$ is multiplicative because $S_{12}$ is a line. The recursion says that the Dirichlet series of $\tau$ has an Euler product, like $\zeta$ and the $L(s,\chi)$ of this book:
\[
\sum_{n\ge1}\frac{\tau(n)}{n^s}=\prod_p\frac1{1-\tau(p)p^{-s}+p^{11-2s}},
\]
and Hecke proved that it extends to an entire function with a functional equation, just as Riemann did for $\zeta$.

\begin{remark}[Ramanujan's third observation, and beyond]
Write $1-\tau(p)X+p^{11}X^2=(1-\alpha_pX)(1-\beta_pX)$. Ramanujan's inequality $|\tau(p)|\le2p^{11/2}$ says exactly that $\alpha_p$ and $\beta_p$ are complex conjugates of modulus $p^{11/2}$: a Riemann hypothesis
for the local factors. Deligne deduced it in 1974 from his proof of the Riemann hypothesis for varieties over finite fields. Write then $\tau(p)=2p^{11/2}\cos\theta_p$. The Sato--Tate conjecture, proved in 2011 by Barnet-Lamb, Geraghty, Harris and Taylor,
says that the angles $\theta_p$ are distributed with density $\frac2\pi\sin^2\theta$ on $[0,\pi]$, so that $x_p=\cos\theta_p$ has density $\frac2\pi\sqrt{1-x^2}$ (Figure~\ref{fig:satotate}). Lehmer asked in 1947 whether $\tau(n)$ is ever $0$;
nobody knows. The same circle of ideas, for modular forms of weight $2$ attached to elliptic curves, is how Wiles proved Fermat's Last Theorem in 1995.
\end{remark}

\begin{figure}[ht]
\centering
\fig{ch12-satotate}
\caption{The normalised values $x_p=\frac{\tau(p)}{2p^{11/2}}$ for the $5133$ primes $p\le50\,000$ (histogram), against the Sato--Tate density $\frac2\pi\sqrt{1-x^2}$ (orange). All the values lie in $[-1,1]$, as Deligne's
theorem says; the largest $|x_p|$ is $0.990$.}
\label{fig:satotate}
\end{figure}

\begin{beauty}
Ramanujan's three observations about the coefficients of one power series turned out to be shadows of three structures: a finite-dimensional space (the $691$), a commuting family of operators acting on it
(multiplicativity), and the geometry of varieties over finite fields (Deligne's bound). The symmetry group of Figure~\ref{fig:tessellation} is the reason for all three. This book began with Gauss counting primes; it ends at the door
of the modern theory, where the primes, the $L$-functions of Chapters~3 to~7, and modular forms are parts of a single landscape. Tome~12 opens that door.
\end{beauty}

\section{Exercises}

\begin{exercise}[$\star$]\label{ex:sigma9}
Prove that $11\sigma_9(n)=21\sigma_5(n)-10\sigma_3(n)+5040\sum_{m=1}^{n-1}\sigma_3(m)\sigma_5(n-m)$ for every $n\ge1$.
\end{exercise}

\begin{exercise}[$\star$]\label{ex:e2}
Let $E_2=1-24\sum_{n\ge1}\sigma_1(n)q^n$. Show that $E_2$ is holomorphic on $\mathbb H$ and $E_2(\tau+1)=E_2(\tau)$, but that $E_2$ is not a modular form of weight $2$.
\end{exercise}

\begin{exercise}[$\star$]\label{ex:zeros}
Show that $E_6(i)=0$ and $E_4(\rho)=0$, where $\rho=e^{2\pi i/3}$. Deduce that $\Delta(i)=\frac{E_4(i)^3}{1728}$.
\end{exercise}

\begin{exercise}[$\star$]\label{ex:tausmall}
Using Theorem~\ref{thm:mordell} and $\tau(2)=-24$, $\tau(3)=252$, compute $\tau(4)$, $\tau(6)$, $\tau(8)$ and $\tau(12)$, and check them against the product. Check Corollary~\ref{cor:delta} for $n=2$ and $n=3$.
\end{exercise}

\begin{exercise}[$\star\star$]\label{ex:eisensteinhecke}
Show that $T_pE_k=(1+p^{k-1})E_k$, and that $T_pT_\ell=T_\ell T_p$ on $M_k$ for distinct primes $p,\ell$.
\end{exercise}

\begin{exercise}[$\star\star$, the signature exercise]\label{ex:heckebound}
(a) Show that $\phi(\tau)=(\operatorname{Im}\tau)^6|\Delta(\tau)|$ is invariant under $\mathrm{SL}_2(\Z)$ and bounded on $\mathbb H$.
(b) Deduce Hecke's bound $\tau(n)=O(n^6)$, by writing $\tau(n)=\int_0^1\Delta(x+iy)e^{-2\pi in(x+iy)}\dd x$ and choosing $y=\frac1n$. (Deligne's theorem gives $O(n^{11/2+\varepsilon})$.)
\end{exercise}

\section{Hints}

\begin{hint}{ex:sigma9}
$E_4E_6=E_{10}$. Divide by $-24$.
\end{hint}
\begin{hint}{ex:e2}
If it were modular, Theorem~\ref{thm:dimension}(a) would apply.
\end{hint}
\begin{hint}{ex:zeros}
$S$ fixes $i$; and $\rho$ is fixed by $ST$, or note $-\frac1\rho=\rho+1$.
\end{hint}
\begin{hint}{ex:tausmall}
$\tau(4)=\tau(2)^2-2^{11}$; $\tau(8)=\tau(2)\tau(4)-2^{11}\tau(2)$.
\end{hint}
\begin{hint}{ex:eisensteinhecke}
Use the $q$-expansion of $T_p$ and the multiplicativity of $\sigma_{k-1}$; for commutation, compute the coefficient of $q^n$ in $T_pT_\ell f$, which is symmetric in $p$ and $\ell$.
\end{hint}
\begin{hint}{ex:heckebound}
(a) $\operatorname{Im}(\gamma\tau)=\frac{\operatorname{Im}\tau}{|c\tau+d|^2}$. On $F$, $\Delta=q\,u(q)$ with $u$ bounded, so $\phi\le(\operatorname{Im}\tau)^6e^{-2\pi\operatorname{Im}\tau}\sup|u|$. (b) $|\Delta(x+iy)|\le Cy^{-6}$.
\end{hint}

\section{Solutions}

\begin{solution}{ex:sigma9}
\idea $M_{10}$ is a line, so the two forms $E_4E_6$ and $E_{10}$ with constant term $1$ must be equal.
\writeup Both have constant term $1$, and $\dim M_{10}=1$, so $E_4E_6=E_{10}$. The coefficient of $q^n$ gives
$-264\sigma_9(n)=240\sigma_3(n)-504\sigma_5(n)-240\cdot504\sum_{m=1}^{n-1}\sigma_3(m)\sigma_5(n-m)$. Divide by $-24$: $11\sigma_9(n)=-10\sigma_3(n)+21\sigma_5(n)+5040\sum\sigma_3(m)\sigma_5(n-m)$. For $n=2$: $11\cdot513=5643$ and
$21\cdot33-10\cdot9+5040=693-90+5040=5643$.
\end{solution}

\begin{solution}{ex:e2}
\idea The proof of Theorem~\ref{thm:dimension} kills every modular form of weight $2$, and $E_2$ has constant term $1$, so it cannot be one. (It fails only by a little: $E_2(-\frac1\tau)=\tau^2E_2(\tau)+\frac{6\tau}{\pi i}$.)
\writeup Since $\sigma_1(n)\le n^2$, the series converges uniformly on $|q|\le r<1$, so $E_2$ is holomorphic in $q$, hence in $\tau$, and is $1$-periodic. If it were a modular form of weight $2$, Theorem~\ref{thm:dimension}(a) would give $E_2=0$,
but its constant term is $1$.
\end{solution}

\begin{solution}{ex:zeros}
\idea A fixed point of a symmetry is a place where the automorphy factor must equal $1$, unless the form vanishes.
\writeup $S(i)=i$, so $E_6(i)=E_6(-\frac1i)=i^6E_6(i)=-E_6(i)$, hence $E_6(i)=0$. For $\rho=e^{2\pi i/3}=-\frac12+\frac{\sqrt3}2i$, $\rho^2+\rho+1=0$ gives $-\frac1\rho=\rho+1$. So
$E_4(\rho)=E_4(\rho+1)=E_4(-\frac1\rho)=\rho^4E_4(\rho)=\rho E_4(\rho)$, and $\rho\ne1$ forces $E_4(\rho)=0$. Then $1728\Delta(i)=E_4(i)^3-E_6(i)^2=E_4(i)^3$; as $\Delta(i)\ne0$, $E_4(i)\ne0$.
\end{solution}

\begin{solution}{ex:tausmall}
\idea Mordell's relations generate all values from the values at primes.
\writeup $\tau(4)=576-2048=-1472$; $\tau(6)=-24\cdot252=-6048$; $\tau(8)=\tau(2)\tau(4)-2^{11}\tau(2)=35328+49152=84480$; $\tau(12)=\tau(4)\tau(3)=-370944$. These match the product, as the notebook checks. For the congruence:
$\sigma_{11}(2)=2049$ and $-24-2049=-2073=-3\cdot691$; $\sigma_{11}(3)=177148$ and $177148-252=176896=256\cdot691$.
\end{solution}

\begin{solution}{ex:eisensteinhecke}
\idea Hecke operators act on $q$-expansions by a rule that only involves the multiplicative structure of the indices; divisor sums are already multiplicative, which makes $E_k$ an eigenform.
\writeup Let $s(n)=\sigma_{k-1}(n)$, with $s(n/p)=0$ if $p\nmid n$. We claim $s(pn)+p^{k-1}s(n/p)=(1+p^{k-1})s(n)$. Write $n=p^rm$ with $p\nmid m$. The divisors of $p^{r+1}m$ are those of $p^rm$ and the $p^{r+1}e$ with $e\mid m$, so
$s(pn)=s(n)+p^{(r+1)(k-1)}s(m)$; similarly, if $r\ge1$, $p^{k-1}s(n/p)=p^{k-1}\bigl(s(n)-p^{r(k-1)}s(m)\bigr)$, and if $r=0$ both $p^{k-1}s(n/p)$ and $p^{k-1}(s(n)-p^{r(k-1)}s(m))$ are $0$. Adding gives the claim. For $n=0$
the constant terms give $(1+p^{k-1})a_0$. So $T_pE_k=(1+p^{k-1})E_k$.

For commutation, write $T_pf=\sum b_nq^n$ with $b_n=a_{pn}+p^{k-1}a_{n/p}$. The coefficient of $q^n$ in $T_\ell T_pf$ is
\[
\begin{aligned}b_{\ell n}+\ell^{k-1}b_{n/\ell}&=a_{p\ell n}+p^{k-1}a_{\ell n/p}\eqbr+\ell^{k-1}a_{pn/\ell}+(p\ell)^{k-1}a_{n/(p\ell)},\end{aligned}
\]
with the convention $a_x=0$ when $x$ is not an integer, which is consistent because $p\ne\ell$: for instance $\frac{pn}\ell$ is an integer exactly when $\ell\mid n$, that is, exactly when $b_{n/\ell}$ is defined.
This expression is symmetric in $p$ and $\ell$, so $T_\ell T_p=T_pT_\ell$.
\end{solution}

\begin{solution}{ex:heckebound}
\idea The weight is exactly what makes $(\operatorname{Im}\tau)^{k/2}|f|$ invariant. An invariant function is controlled by its values on $F$, where a cusp form decays exponentially. Then the Fourier coefficient is an
integral along a horizontal line, and the best line is at height $\frac1n$.
\writeup (a) $\phi(\gamma\tau)=\frac{(\operatorname{Im}\tau)^6}{|c\tau+d|^{12}}|c\tau+d|^{12}|\Delta(\tau)|=\phi(\tau)$. By Theorem~\ref{thm:fundamental}, $\sup_{\mathbb H}\phi=\sup_F\phi$. On $F$, $|q|\le e^{-\pi\sqrt3}$ and $\Delta=q\,u(q)$ with
$u$ continuous on the closed disc $|q|\le e^{-\pi\sqrt3}$, so $\phi(\tau)\le\max|u|\cdot y^6e^{-2\pi y}$ with $y=\operatorname{Im}\tau$, which is bounded. So $\phi\le C$ on $\mathbb H$.

(b) For fixed $y>0$, the series
\[
\Delta(x+iy)=\sum_m\tau(m)e^{-2\pi my}e^{2\pi imx}
\]
converges uniformly in $x$, so integrating term by term over $[0,1]$,
\[
\tau(n)e^{-2\pi ny}=\int_0^1\Delta(x+iy)e^{-2\pi inx}\dd x.
\]
Hence $|\tau(n)|\le e^{2\pi ny}\max_x|\Delta(x+iy)|$, which is at most $e^{2\pi ny}Cy^{-6}$. With $y=\frac1n$: $|\tau(n)|\le e^{2\pi}Cn^6$.
\end{solution}

\section*{Sources and further reading}

\begin{itemize}
\item S.~Ramanujan, \emph{On certain arithmetical functions}, Transactions of the Cambridge Philosophical Society 22 (1916), 159--184.
\item L.~J.~Mordell, \emph{On Mr Ramanujan's empirical expansions of modular functions}, Proceedings of the Cambridge Philosophical Society 19 (1917), 117--124.
\item E.~Hecke, \emph{Über Modulfunktionen und die Dirichletschen Reihen mit Eulerscher Produktentwicklung I, II}, Mathematische Annalen 114 (1937), 1--28 and 316--351.
\item C.~L.~Siegel, \emph{A simple proof of $\eta(-1/\tau)=\eta(\tau)\sqrt{\tau/i}$}, Mathematika 1 (1954), 4.
\item P.~Deligne, \emph{La conjecture de Weil I}, Publications Mathématiques de l'IHÉS 43 (1974), 273--307.
\item J.-P.~Serre, \emph{A Course in Arithmetic}, Springer GTM 7, 1973, Chapter~VII. The classic short introduction.
\item T.~M.~Apostol, \emph{Modular Functions and Dirichlet Series in Number Theory}, 2nd ed., Springer GTM 41, 1990, Chapters~1--3 and~6. Siegel's proof and Hecke operators.
\item D.~Zagier, \emph{Elliptic modular forms and their applications}, in \emph{The 1-2-3 of Modular Forms}, Springer, 2008, 1--103. A delightful survey.
\item F.~Diamond and J.~Shurman, \emph{A First Course in Modular Forms}, Springer GTM 228, 2005.
\item T.~Barnet-Lamb, D.~Geraghty, M.~Harris and R.~Taylor, \emph{A family of Calabi--Yau varieties and potential automorphy II}, Publications of the RIMS 47 (2011), 29--98. The Sato--Tate conjecture.
\end{itemize}
